\section*{Chapter \ref{ch:b3:probability} --- Probability:
Foundations and the Law of Large Numbers}

\begin{solution}{exo:b3:probability:1}
(a) For $u \in \intoo01$: $\P(F(X) \leq u) = \P(X \leq
F^{-1}(u)) = F(F^{-1}(u)) = u$ (continuity and strict
monotonicity make $F$ a bijection onto $\intoo01$ with
$\{F(X) \leq u\} = \{X \leq F^{-1}(u)\}$): $F(X)$ is uniform.
Conversely $\P(G(U) \leq t) = \P(U \leq F(t)) = F(t)$: to
simulate a law, apply the inverse distribution function to a
uniform sample.

(b) $Y = X^2$, $X$ uniform on $\intcc{-1}1$: for $t \in
\intcc01$, $F_Y(t) = \P(-\sqrt t \leq X \leq \sqrt t) = \sqrt
t$: density $\frac1{2\sqrt t}\mathbf 1_{\intoo01}$. And
$\P\bigl(-\frac1\lambda\ln U \leq t\bigr) = \P(U \geq
\eu^{-\lambda t}) = 1 - \eu^{-\lambda t}$: the exponential
$\mathcal E(\lambda)$ --- inversion in action.
\end{solution}

\begin{solution}{exo:b3:probability:2}
(a) Poisson: $\E X = \sum_{k\geq0}k\,\eu^{-\lambda}
\frac{\lambda^k}{k!} = \lambda$, $\E[X(X-1)] = \lambda^2$, so
$\V = \lambda^2 + \lambda - \lambda^2 = \lambda$. Geometric
($\P(X = k) = p(1-p)^{k-1}$): $\E X = \frac1p$, $\V =
\frac{1-p}{p^2}$ (differentiate the geometric series twice).

(b) $G(t) = \P(T > t)$ is nonincreasing with $G(0^+)\dots$
$G \colon \intco0\infty \to \intoc01$; memorylessness reads
$G(t + s) = G(t)G(s)$. Then $G(n t) = G(t)^n$ and $G(t/n) =
G(t)^{1/n}$: $G(q) = G(1)^q$ for rational $q \geq 0$; writing
$G(1) = \eu^{-\lambda}$ ($\in \intoo01$: $G(1) = 1$ would
force $G \equiv 1$, impossible for a finite random variable;
$G(1) = 0$ is excluded by hypothesis) and squeezing an
arbitrary $t$ between rationals (monotonicity): $G(t) =
\eu^{-\lambda t}$ --- the exponential law. The converse is a
computation.
\end{solution}

\begin{solution}{exo:b3:probability:3}
(a) $\sigma(f_i(X_i)) = f_i(X_i)^{-1}(\mathcal B) \subseteq
X_i^{-1}(\mathcal B) = \sigma(X_i)$ ($f_i$ Borel), and
sub-$\sigma$-algebras of independent $\sigma$-algebras are
independent (the defining identity holds a fortiori).

(b) $\sigma(A_i) = \{\varnothing, A_i, A_i^c, \Omega\} =
\sigma(A_i^c) = \sigma(\mathbf 1_{A_i})$: all three
statements assert independence of the same
$\sigma$-algebras. (That factorization over the $A_i$
propagates to complements is the $\lambda$-system argument
inside \cref{def:b3:probability:independence}'s equivalence
--- or direct inclusion-exclusion.)

(c) $\P(A) = \P(B) = \P(C) = \frac12$; $A\cap B = A\cap C =
B\cap C$ on pairs: each intersection is ``both heads'' or
analogous, of probability $\frac14$: pairwise independent.
But $\P(A\cap B\cap C) = \P(\text{HH}) = \frac14 \neq
\frac18$: not independent --- $C$ is determined by $A$ and
$B$.
\end{solution}

\begin{solution}{exo:b3:probability:4}
(a) Let the text $T$ have length $L$ and $q = a^{-L}$ ($a$
the alphabet size). The events $E_k = \{$positions $kL+1,
\dots, (k+1)L$ spell $T\}$ are independent (disjoint blocks
of i.i.d.\ letters), each of probability $q > 0$: $\sum\P(E_k)
= \infty$, and Borel--Cantelli (2) gives infinitely many
occurrences a.s.

(b) Upper: $\P\bigl(R_n \geq (1+\varepsilon)\log_2n\bigr)
\leq 2^{-(1+\varepsilon)\log_2n} = n^{-(1+\varepsilon)}$,
summable: by Borel--Cantelli (1), a.s.\ only finitely many
such $n$. Lower: pack disjoint blocks --- the $j$-th of
length $\ell_j = \lceil\log_2s_j\rceil$ starting at $s_j =
\sum_{i<j}\ell_i$; the events ``block $j$ is all ones'' are
independent with probability $2^{-\ell_j} \asymp
\frac1{s_j} \asymp \frac1{j\log_2 j}$, whose sum diverges:
Borel--Cantelli (2) gives infinitely many all-ones blocks,
i.e.\ $R_{s_j} \geq \log_2 s_j$ infinitely often. Together:
the maximal run length in the first $n$ digits is
$(1 + o(1))\log_2n$ a.s.
\end{solution}

\begin{solution}{exo:b3:probability:5}
(a) $R = \bigl(\limsup\abs{X_n}^{1/n}\bigr)^{-1}$ is
unchanged if finitely many $X_n$ are modified: for every $N$,
$R$ is $\sigma(X_N, X_{N+1}, \dots)$-measurable, i.e.\
tail-measurable. Then each event $\{R \leq c\}$ has
probability $0$ or $1$ (\cref{thm:b3:probability:zeroone}),
so the distribution function of $R$ takes only the values
$0, 1$: it jumps at a single point $c_0 \in
\intcc0{+\infty}$, and $R = c_0$ a.s.

(b) Convergence of $\sum X_n$ and of $\frac{S_n}n$ are
insensitive to changing finitely many terms (for the second:
the modified terms contribute $O(1/n) \to 0$): tail events;
zero--one law.

(c) $\{X_1 > 0\}$ depends on $X_1$: for i.i.d.\ signs
($\P(X_1 = \pm1) = \frac12$), its probability is $\frac12
\notin \{0,1\}$ --- no contradiction, it is not a tail
event.
\end{solution}

\begin{solution}{exo:b3:probability:6}
Work on $(\intcc01, \lambda)$.
(a) The typewriter $\mathbf 1_{I_n}$
(\cref{exo:b3:lp:3}): $\norm{X_n}_p^p = \lambda(I_n) \to 0$
(all $p < \infty$), hence also in probability; at every
$\omega$ the values $0$ and $1$ both recur: no pointwise
convergence anywhere.
(b) $X_n = n\mathbf 1_{\intoo0{1/n}} \to 0$ off $0$, but
$\norm{X_n}_p \geq n^{1 - 1/p} \geq 1$.
(c) $X_n = \sqrt n\,\mathbf 1_{\intoo0{1/n}}$: $\E\abs{X_n} =
n^{-1/2} \to 0$, $\E X_n^2 = 1$.
(d) From any subsequence extract (convergence in
probability) a further subsequence converging a.s.\
(\cref{prop:b3:probability:modes}(c)); dominated convergence
gives $L^1$ convergence along it, with the \emph{same} limit
$X$. Thus every subsequence of the numerical sequence
$\E\abs{X_n - X}$ has a subsubsequence tending to $0$: the
whole sequence tends to $0$.
\end{solution}

\begin{solution}{exo:b3:probability:7}
(a) $\hat p_n = \frac{S_n}n$ with $S_n$ binomial: $\V(\hat
p_n) = \frac{p(1-p)}n \leq \frac1{4n}$, and Chebyshev
(\cref{prop:b3:probability:markov}) gives the bound.
(b) Solve $\frac1{4n(0.03)^2} \leq 0.05$: $n \geq
\frac{1}{4\cdot0.0009\cdot0.05} \approx 5556$. The central
limit theorem will justify $n \approx 1070$ for the same
guarantee: Chebyshev pays for its generality with a factor
$\approx 5$.
\end{solution}

\begin{solution}{exo:b3:probability:8}
(a) If $S_n \sim \mathcal B(n, x)$, the transfer theorem
gives $\E\bigl[f(\frac{S_n}n)\bigr] =
\sum_k\binom nkx^k(1-x)^{n-k}f(\frac kn) = B_nf(x)$.

(b)--(c) Let $\omega = \omega_f$ be the modulus of continuity
($\abs{f(u) - f(v)} \leq \omega(\abs{u - v})$, and $\omega(c
\delta) \leq (1 + c)\,\omega(\delta)$ by chaining steps).
Then, for any $\delta > 0$,
\[
\abs{f(u) - f(x)} \leq \Bigl(1 + \frac{(u -
x)^2}{\delta^2}\Bigr)\omega(\delta)
\]
(if $\abs{u - x} \leq \delta$, clear; otherwise
$\omega(\abs{u-x}) \leq (1 + \frac{\abs{u-x}}\delta)
\omega(\delta) \leq (1 + \frac{(u-x)^2}{\delta^2})
\omega(\delta)$). Take expectations at $u = \frac{S_n}n$:
\[
\abs{B_nf(x) - f(x)} \leq
\Bigl(1 + \frac{\V(S_n/n)}{\delta^2}\Bigr)\omega(\delta)
\leq \Bigl(1 + \frac{1}{4n\delta^2}\Bigr)\omega(\delta) ;
\]
with $\delta = n^{-1/2}$: $\norm{B_nf - f}_\infty \leq
\frac54\,\omega\bigl(n^{-1/2}\bigr) \leq
\frac32\,\omega\bigl(n^{-1/2}\bigr) \to 0$ (uniform
continuity on the compact): a probabilistic Weierstrass
theorem, with an explicit and uniform rate.
\end{solution}

\begin{solution}{exo:b3:probability:9}
(a) After $k - 1$ types are collected, each draw is new with
probability $p_k = \frac{n-k+1}n$: $\tau_k$ is geometric
$(p_k)$, and the $\tau_k$ are independent (the draws are).
Sums: $\E T_n = \sum_k\frac n{n-k+1} = nH_n \sim n\ln n$;
$\V(T_n) = \sum\frac{1 - p_k}{p_k^2} \leq
n^2\sum_{j=1}^n\frac1{j^2} \leq \frac{\pi^2}6n^2$.

(b) Chebyshev: $\P\bigl(\abs{T_n - nH_n} \geq \varepsilon
n\ln n\bigr) \leq \frac{\pi^2n^2/6}{\varepsilon^2n^2\ln^2n}
\to 0$, and $\frac{nH_n}{n\ln n} \to 1$: $\frac{T_n}{n\ln n}
\to 1$ in probability.

(c) Union bound: $T_n > t$ means some type is unseen after
$\lceil t\rceil$ draws, so $\P(T_n > t) \leq n(1 -
\frac1n)^{t} \leq n\,\eu^{-t/n}$; at $t = \beta n\ln n$:
$\leq n^{1 - \beta}$. For $\beta > 1$, $\sum_m
2^{m(1-\beta)} < \infty$: Borel--Cantelli gives, along $n =
2^m$, a.s.\ $T_n \leq \beta n\ln n$ eventually --- in
particular for every $\beta > 2$ as stated (any $\beta > 1$
works along the subsequence).
\end{solution}

\begin{solution}{exo:b3:probability:10}
(a) The even-indexed digits $(b_{2k})_k$ are i.i.d.\ fair
bits (a subfamily of the independent digit family), so $U =
\sum_kb_{2k}2^{-k}$ gives every dyadic interval its correct
probability (as in
\cref{thm:b3:probability:existence}): uniform; likewise $V$;
and $(U, V)$ depend on disjoint digit blocks: independent
(factorization on dyadic rectangles, then Dynkin).

(b) $\Phi(\omega) = (U(\omega), V(\omega))$ is measurable
with $\Phi_*\lambda = \lambda\otimes\lambda = \lambda_2$
(agreement on dyadic rectangles + uniqueness). Interleaving
digits defines an inverse defined off the (null) set of
dyadic rationals in either factor: a measure-preserving
bijection between full-measure subsets of $\intcc01$ and
$\intcc01^2$. Contrast with Peano
(\cref{pb:b3:topology:1}): continuity forced surjectivity
without injectivity; dropping continuity for mere
measurability buys a measure-isomorphism --- dimension is
invisible to measure theory, visible to topology.
\end{solution}

\begin{solution}{pb:b3:probability:1}
\textbf{1.} $X_n^{\pm}$ are Borel functions of $X_n$: they
remain pairwise independent
(\cref{exo:b3:probability:3}(a)) and identically
distributed, integrable, with $\E X_1 = \E X_1^+ - \E
X_1^-$. If the theorem holds for nonnegative variables,
apply it to both halves and subtract:
$\frac{S_n}n = \frac{S_n^+}n - \frac{S_n^-}n \to \E X_1^+ -
\E X_1^- = m$ a.s.

\textbf{2.} $\P(X_n \neq Y_n) = \P(X_n > n) = \P(X_1 > n)$
(identical laws), and $\sum_n\P(X_1 > n) \leq \sum_n\P(X_1
\geq n) \leq \E X_1 < \infty$
(\cref{exo:b3:product:3}(a)). Borel--Cantelli (1): a.s.\
$X_n = Y_n$ for all large $n$, so $S_n - S_n^*$ is
eventually constant in $n$: $\frac{S_n - S_n^*}n \to 0$
a.s., and the two normalized sums share their asymptotic
behavior.

\textbf{3.} $X_1\mathbf 1_{X_1 \leq n} \nearrow X_1$: MCT
gives $\E Y_n \to m$; Ces\`aro means of a convergent
sequence converge to the same limit: $\frac{\E S_n^*}n =
\frac1n\sum_{k\leq n}\E Y_k \to m$. Hence it suffices to
prove $\frac{S^*_n - \E S^*_n}{n} \to 0$ a.s.

\textbf{4.} $\V(Y_n) \leq \E Y_n^2 = \E[X_1^2\mathbf
1_{X_1\leq n}]$. By Tonelli for series,
\[
\sum_n\frac{\E[X_1^2\mathbf 1_{X_1\leq n}]}{n^2}
= \E\Bigl[X_1^2\!\!\sum_{n \geq \max(X_1, 1)}\!\frac1{n^2}
\Bigr]
\leq \E\Bigl[X_1^2\cdot\frac{4}{\max(X_1,1)}\Bigr]
\leq 4\,\E[X_1] < \infty,
\]
using $\sum_{n\geq x}n^{-2} \leq \frac4x$ for $x \geq 1$
(for $x \geq 2$: $\leq \frac1{x-1} \leq \frac2x$; for $1
\leq x < 2$: $\leq \frac{\pi^2}6 \leq \frac4x$ since
$\frac4x > 2$), and $X_1^2/\max(X_1, 1) \leq X_1$ in both
cases $X_1 \gtrless 1$.

\textbf{5.} Pairwise independence gives $\E[(Y_i - \E
Y_i)(Y_j - \E Y_j)] = 0$ for $i \neq j$ (the product
formula for two variables), so variances add: $\V(S^*_k) =
\sum_{n\leq k}\V(Y_n)$. Chebyshev on each $k_j$ and
summing:
\[
\sum_j\P\Bigl(\abs{S^*_{k_j} - \E S^*_{k_j}} \geq
\varepsilon k_j\Bigr)
\leq \frac1{\varepsilon^2}\sum_j\frac1{k_j^2}\sum_{n\leq
k_j}\V(Y_n)
= \frac1{\varepsilon^2}\sum_n\V(Y_n)\!\!\sum_{j : k_j\geq
n}\!\frac1{k_j^2}
\]
(Tonelli for the nonnegative double series).

\textbf{6.} $k_j = \lfloor\alpha^j\rfloor \geq
\frac{\alpha^j}2$ (valid once $\alpha^j \geq 1$, i.e.\ all
$j \geq 0$: $\lfloor x\rfloor \geq \frac x2$ for $x \geq
1$). Hence
\[
\sum_{j : k_j \geq n}\frac1{k_j^2}
\leq 4\sum_{j : \alpha^j \geq n}\alpha^{-2j}
\leq \frac{4}{1 - \alpha^{-2}}\cdot\frac1{n^2}
= \frac{C_\alpha}{n^2},
\]
(geometric series from the first $j$ with $\alpha^j \geq
n$). Combining with questions 4--5, the double sum is
finite; Borel--Cantelli (1), applied for each rational
$\varepsilon$ and intersected, gives
$\frac{S^*_{k_j} - \E S^*_{k_j}}{k_j} \to 0$ a.s., and with
question 3: $\frac{S^*_{k_j}}{k_j} \to m$ a.s.

\textbf{7.} $Y_n \geq 0$ makes $n \mapsto S^*_n$
nondecreasing: for $k_j \leq n \leq k_{j+1}$,
\[
\frac{S^*_{k_j}}{k_{j+1}} \leq \frac{S^*_n}{n} \leq
\frac{S^*_{k_{j+1}}}{k_j},
\]
which is the displayed sandwich after inserting
$\frac{k_j}{k_{j+1}}$ and $\frac{k_{j+1}}{k_j}$. Since
$\frac{k_{j+1}}{k_j} \to \alpha$, question 6 gives a.s.
\[
\frac m\alpha \leq \liminf_n\frac{S^*_n}n \leq
\limsup_n\frac{S^*_n}n \leq \alpha m .
\]

\textbf{8.} Apply question 7 for $\alpha = 1 + \frac1p$, $p
\in \N^*$: countably many a.s.\ events; on their
intersection, letting $p \to \infty$: $\lim\frac{S^*_n}n =
m$ a.s. With questions 1--3, $\frac{S_n}n \to \E X_1$
a.s.: the strong law of large numbers, under pairwise
independence.

\textbf{9.} Independence-type hypotheses appeared three
times: (i) additivity of variances (question 5) ---
pairwise suffices; (ii) identical distribution, in the
truncation sums (question 2) and the mean computation
(question 3) --- no independence at all; (iii)
Borel--Cantelli (1) (questions 2 and 6) --- valid without
any independence. Full mutual independence was never
invoked: Etemadi's observation.

\textbf{10.} Fix a base $b$ and a digit $r$. The base-$b$
digits $(d_k)$ of a uniform $\omega$ are i.i.d.\ uniform on
$\{0, \dots, b-1\}$ (each digit vector value occupies an
interval of length $b^{-m}$: the argument of
\cref{thm:b3:probability:existence} verbatim). The strong
law applied to the i.i.d.\ bounded variables $\mathbf
1_{d_k = r}$ gives: a.s., the frequency of the digit $r$
tends to $\frac1b$. Intersecting over the countably many
pairs $(b, r)$: almost every number is
\emph{simply normal in every base}. An explicit non-normal
number: $x = 0.100100100\ldots_2$ (frequency of ones
$\frac13 \neq \frac12$). The contrast is humbling: almost
all numbers are normal, yet for $\sqrt2$, $\eu$, or $\pi$
normality remains unproved --- measure theory counts
without exhibiting.

\textbf{11.} By \cref{exo:b3:probability:10} iterated, a
single uniform variable yields a sequence of i.i.d.\
uniform \emph{vectors} $U_k$ on $\intcc01^d$ (split the
digit set of each $U_n$ of
\cref{thm:b3:probability:existence} into $d$ subfamilies).
For $g \in L^1(\intcc01^d)$, the variables $g(U_k)$ are
i.i.d.\ integrable with mean $\int g\,\dd\lambda_d$
(transfer): the strong law gives
\[
\frac1n\sum_{k=1}^ng(U_k)
\xrightarrow[n\to\infty]{\text{a.s.}}
\int_{\intcc01^d}g\,\dd\lambda_d :
\]
Monte Carlo integration converges almost surely, in every
dimension --- the error size is the business of the central
limit theorem (\cref{ch:b3:clt}).
\end{solution}
